\section{Model Details}
\label{supp:model}

This appendix expands the model description in Sec.~\ref{sec:pretraining} and
Sec.~\ref{sec:posttraining}. \molmoactt is built in two architectural stages.
First, \molmoacttpre adapts the \molmoer vision-language backbone into a
discrete autoregressive robot policy. It keeps the Molmo2 visual and text token
interface, but adds robot-specific state, setup/control, and action-output
tokens so that robot trajectories can be trained with the same next-token
objective as vision-language data. Second, post-training attaches a continuous
DiT-style action expert to this already robot-aware backbone. The expert uses
per-layer KV conditioning on the VLM and learns a flow-matching velocity field
over continuous robot action chunks. \molmothink keeps the same architecture and
adds an autoregressive depth-token prefix before action prediction.

\subsection{\molmoactt Backbone}
\label{supp:model-backbone}

\paragraph{Backbone initialization and visual inputs.}
\molmoactt initializes from \molmoer, a Molmo2-based VLM specialized for
embodied and spatial reasoning. The backbone therefore follows the Molmo2
architecture~\citep{clark2026molmo2openweightsdata}: visual inputs are encoded
by a SigLIP2 vision transformer, pooled and projected by a vision-language
connector, and consumed by an autoregressive language model together with text.
The important MolmoAct2-specific choice is how this interface is used for robot
control. Robot examples can include multiple camera views, proprioceptive state
tokens, setup/control descriptors, and action targets, so sequence length is a
central constraint. For the robot-policy stages in Sec.~\ref{sec:pretraining}
and Sec.~\ref{sec:posttraining}, each camera observation is represented as a
single resized crop rather than a tiled high-resolution image. Video examples
are sampled at up to 2 FPS and capped to at most 8 frames. Multi-camera robot
episodes are serialized as multiple image inputs, and the camera order is
randomized at the episode level during pre-training.

\paragraph{Vision encoder and connector.}
Each crop is encoded by a SigLIP2 vision transformer~\citep{tschannen2025siglip}.
The connector follows the Molmo2 design. It reads patch features from the
third-to-last and ninth-from-last ViT layers, then pools local windows with a
multi-head attention layer. For images, \(2\times2\) patch windows are pooled
into one visual token using the mean patch feature as the query. For videos, a
\(3\times3\) pooling window is used to further reduce the number of frame
tokens. The pooled features are projected into the language-model embedding
space with a shared MLP. This produces a sequence of visual tokens that can be
interleaved with text tokens.

\paragraph{Language model input.}
The language model receives visual tokens together with the robot instruction
and robot-specific text. Multi-camera observations are identified by image-index
text, and video frames use the same frame markers as Molmo2. Since MolmoAct2
robot observations use single resized crops, the model does not rely on the
multi-crop column-token machinery for robot control. Visual tokens are allowed
to forward-attend to one another, including across different frames or camera
views, so the backbone can reason over the full visual context before producing
state-conditioned robot outputs.

\paragraph{Added tokens.}
The robot interface keeps the model in the same autoregressive format as the
base VLM by adding robot-specific tokens to the backbone vocabulary. A robot
example contains visual observations, a language instruction, optional setup
and control descriptors, discrete state tokens, and an output marker that tells
the model which response type to produce. The main markers are
\texttt{\detokenize{<setup_start>}}, \texttt{\detokenize{<setup_end>}},
\texttt{\detokenize{<control_start>}}, \texttt{\detokenize{<control_end>}},
\texttt{\detokenize{<state_start>}}, \texttt{\detokenize{<state_end>}}, and
\texttt{\detokenize{<action_output>}}. We also add indexed token families for
discrete robot quantities: action tokens
\texttt{\detokenize{<action_0>}}--\texttt{\detokenize{<action_2047>}} and
state tokens
\texttt{\detokenize{<state_0>}}--\texttt{\detokenize{<state_255>}}. The
depth tokens are added only for \molmothink post-training: the model uses
\texttt{\detokenize{<depth_output>}}, \texttt{\detokenize{<depth_start>}}, and
\texttt{\detokenize{<depth_end>}}, together with depth-code tokens
\texttt{\detokenize{<depth_0>}}--\texttt{\detokenize{<depth_127>}}.

Setup and control descriptors make embodiment and action semantics explicit in
the prompt. For example, a bimanual YAM prompt can indicate that the robot is a
pair of YAM arms and that the control target is an absolute joint pose, while a
DROID prompt can indicate a Franka arm with delta end-effector control. This
lets the same tokenizer and backbone support different robots without forcing
all datasets into a single physical coordinate convention.

Continuous robot states are normalized using dataset statistics and discretized
into one of 256 state tokens per scalar dimension. The resulting state-token
sequence is appended to the prompt before the action output marker. Continuous
actions are normalized, padded to a maximum width of 32 dimensions, and encoded
by \molmoactfast into discrete action tokens with a 2048-token action
vocabulary. During pre-training, these action tokens are the only robot target.
During post-training and fine-tuning, the same example also provides a
continuous action chunk for the action expert. The discrete action-token span is
masked from the expert conditioning path, so the continuous policy cannot
condition on the ground-truth action tokens it is trained to predict.

% ---------- MODEL ARCHITECTURE ---------------------------------------------
\begin{table}[t]
\newcommand\modelvalue[1]{\cellcolor{olive!5} #1}
\newcommand\navalue{--}
\centering
\footnotesize
\setlength\tabcolsep{6pt}
\renewcommand\arraystretch{1.1}
\begin{tabular}{@{}lcccc@{}}
\toprule
\textbf{Hyperparameter} &
\textbf{\makecell{Image\\Encoder}} &
\textbf{\makecell{V/L\\Connector}} &
\textbf{LLM} &
\textbf{\makecell{Action\\Expert}} \\
\midrule
Params & \modelvalue{380M} & \modelvalue{57M} & \modelvalue{4.0B} & \modelvalue{621M} \\
Dim & \modelvalue{1152} & \modelvalue{1152} & \modelvalue{2560} & \modelvalue{768} \\
MLP Dim & \modelvalue{4304} & \modelvalue{9728} & \modelvalue{9728} & \modelvalue{3072} \\
Act. & \modelvalue{GELU} & \modelvalue{SwiGLU} & \modelvalue{SwiGLU} & \modelvalue{SwiGLU} \\
Heads & \modelvalue{16} & \modelvalue{16} & \modelvalue{32} & \modelvalue{8} \\
KV Heads & \modelvalue{16} & \navalue & \modelvalue{8} & \navalue \\
Layers & \modelvalue{27} & \navalue & \modelvalue{36} & \modelvalue{36} \\
Dropout & \modelvalue{0.0} & \modelvalue{0.0} & \modelvalue{0.1} & \modelvalue{0.0} \\
\midrule
Image Size & \modelvalue{\(384{\times}384\)} & \navalue & \navalue & \navalue \\
Patch Size & \modelvalue{14} & \navalue & \navalue & \navalue \\
Image Pool Size & \navalue & \modelvalue{\(2{\times}2\)} & \navalue & \navalue \\
Video Pool Size & \navalue & \modelvalue{\(3{\times}3\)} & \navalue & \navalue \\
Pool Dim & \navalue & \modelvalue{1152} & \navalue & \navalue \\
Pool Heads & \navalue & \modelvalue{16} & \navalue & \navalue \\
Embed & \navalue & \navalue & \modelvalue{151936} & \navalue \\
Theta & \navalue & \navalue & \modelvalue{1M} & \navalue \\
\midrule
Max Horizon & \navalue & \navalue & \navalue & \modelvalue{30} \\
Max Action Dim & \navalue & \navalue & \navalue & \modelvalue{32} \\
Time Embed Dim & \navalue & \navalue & \navalue & \modelvalue{256} \\
Conditioning & \navalue & \navalue & \navalue & \modelvalue{Per-layer KV} \\
\bottomrule
\end{tabular}
\caption{\textbf{\molmoactt architecture hyperparameters.} The image encoder,
connector, and LLM columns follow the 4B \molmoer backbone used by
\molmoactt. The action-expert column describes the MolmoAct2 continuous
flow-matching module used for continuous action prediction.}
\label{tab:hyperparams_model}
\end{table}


\subsection{Continuous Action Expert}
\label{supp:model-action-expert}

\paragraph{Input and output.}
The action expert is a noncausal transformer over a short action chunk. It maps
a noisy normalized trajectory
\(x_t \in \mathbb{R}^{H \times D_{\max}}\) to a velocity prediction of the same
shape, where \(H\) is the action horizon and \(D_{\max}=32\) is the shared
maximum action width. The released \molmoactt checkpoint uses \(H=30\) for
post-training, while the \libero fine-tuned checkpoints use \(H=10\). The
corresponding padding masks remove padded action steps and padded action
dimensions from both the noisy input and the loss.

\paragraph{Flow-matching objective.}
Let \(a\) be a normalized target action chunk and let
\(\epsilon \sim \mathcal{N}(0,I)\). For a sampled time \(t \in [0,1]\), we
construct
\begin{equation}
    x_t = (1-t)\epsilon + t a,
    \qquad
    u^\star = a - \epsilon .
\end{equation}
The action expert \(f_\theta\) predicts \(u^\star\) from the noisy trajectory,
the time embedding, and the VLM context \(c\):
\begin{equation}
    \mathcal{L}_{\mathrm{flow}}
    =
    \mathbb{E}_{a,\epsilon,t}
    \left[
    \left\|
    m \odot \left(f_\theta(x_t,t,c) - u^\star\right)
    \right\|_2^2
    \right],
\end{equation}
where \(m\) masks padded action steps and padded action dimensions. In
post-training, we evaluate multiple sampled flow-matching times for each robot
action chunk while reusing the same VLM context. At inference, we initialize
the trajectory from Gaussian noise and integrate the learned velocity field for
a fixed number of Euler steps:
\begin{equation}
    z_{i+1} = z_i + \Delta t \, f_\theta(z_i,t_i,c),
    \qquad
    t_i = \frac{i}{N},
    \qquad
    \Delta t = \frac{1}{N}.
\end{equation}
The released checkpoints use \(N=10\) inference steps. The final normalized
trajectory is sliced to the embodiment-specific action width and unnormalized
with the corresponding dataset statistics.

\paragraph{Expert block.}
The action expert has one transformer block for each VLM layer, giving
\(L=36\) expert blocks. A noisy action chunk is first projected from 32
continuous dimensions to the expert hidden width. Each block contains
bidirectional self-attention over the action chunk, cross-attention to the VLM
context, and an MLP. A sinusoidal time embedding is passed through a small MLP,
then used to produce DiT-style shift, scale, and gate parameters for all three
residual branches. Schematically, block \(\ell\) computes
\begin{align}
    h'_\ell
        &= h_\ell
        + g^{\mathrm{sa}}_\ell
        \operatorname{SA}\!\left(
        \operatorname{AdaRMS}^{\mathrm{sa}}_\ell(h_\ell,t)
        \right), \\
    \bar{h}_\ell
        &= h'_\ell
        + g^{\mathrm{ca}}_\ell
        \operatorname{CA}\!\left(
        \operatorname{AdaRMS}^{\mathrm{ca}}_\ell(h'_\ell,t),
        \tilde{K}_\ell,
        \tilde{V}_\ell
        \right), \\
    h_{\ell+1}
        &= \bar{h}_\ell
        + g^{\mathrm{ff}}_\ell
        \operatorname{MLP}\!\left(
        \operatorname{AdaRMS}^{\mathrm{ff}}_\ell(\bar{h}_\ell,t)
        \right).
\end{align}
Here, \(\operatorname{AdaRMS}\) denotes RMS normalization followed by the
time-dependent affine modulation. The self-attention and cross-attention layers
use query-key normalization. The expert uses rotary position embeddings for the
action sequence, which gives the denoising transformer an explicit ordering
over the predicted trajectory steps. After the final expert block, a modulated
RMS normalization layer and a linear projection map the hidden states back to
32 action dimensions. The final projection is initialized to produce near-zero
velocity predictions at the beginning of post-training, which makes the newly
attached expert easier to optimize.

\paragraph{Per-layer KV conditioning.}
The action expert receives VLM context through per-layer KV conditioning. For
VLM layer \(\ell\), let \((K^{\mathrm{vlm}}_\ell,V^{\mathrm{vlm}}_\ell)\) be
the key and value tensors produced by that layer's self-attention over the
prompt, images, state tokens, and non-target robot text. The expert uses learned
adapter projections \(P_K\) and \(P_V\) to map these tensors into the expert
cross-attention width:
\begin{equation}
    \tilde{K}_\ell =
    \operatorname{reshape}\!\left(P_K K^{\mathrm{vlm}}_\ell\right),
    \qquad
    \tilde{V}_\ell =
    \operatorname{reshape}\!\left(P_V V^{\mathrm{vlm}}_\ell\right).
\end{equation}
Expert block \(\ell\) then cross-attends to
\((\tilde{K}_\ell,\tilde{V}_\ell)\):
\begin{equation}
    \operatorname{CA}(Q_\ell,\tilde{K}_\ell,\tilde{V}_\ell)
    =
    \operatorname{softmax}\!\left(
        \frac{Q_\ell \tilde{K}_\ell^\top}{\sqrt{d_h}}
    \right)\tilde{V}_\ell .
\end{equation}
This gives each expert block direct access to the attention state at the same
depth in the VLM. In the released architecture, the VLM has 8 KV heads with
head dimension 128, so each token contributes a 1024-dimensional key and value
state. The adapter projections map this state to the action expert width of
768, which is reshaped into 8 expert attention heads of width 96. In
post-training, this conditioning path is detached from the VLM for the flow
loss. The language-model loss still updates the VLM, while the flow loss trains
the action expert and the VLM-to-expert adapter projections.

\subsection{Adaptive-Depth Extension}
\label{supp:model-depth}

\molmothink keeps the same backbone and action expert, but adds an
autoregressive depth-token interface before continuous action prediction. A
depth response is requested with \texttt{\detokenize{<depth_output>}} and
serialized between \texttt{\detokenize{<depth_start>}} and
\texttt{\detokenize{<depth_end>}}. The released depth checkpoints predict a
\(10\times10\) depth-token grid, giving 100 depth positions, and each position
takes one of 128 learned depth-code values. For depth-only examples, the
assistant response begins with
\texttt{\detokenize{<depth_output>}} and the model predicts depth tokens. For
depth-and-action examples, the response begins with
\texttt{\detokenize{<depth_output><action_output>}}. The depth tokens are
generated autoregressively, then included in the context used by the action
expert.

The \libero depth fine-tuned model also includes a learned depth gate on the
expert conditioning path. For action-expert layer \(\ell\), let \(M_t=1\)
denote positions belonging to the depth-output trigger, depth delimiters, or
depth-code tokens, and let \(A_t\) denote valid context positions. The gate is
computed from the non-depth context of the corresponding VLM layer,
\begin{equation}
    c_\ell
    =
    \frac{
        \sum_t A_t(1-M_t)V^{\mathrm{vlm}}_{\ell,t}
    }{
        \sum_t A_t(1-M_t)
    },
    \qquad
    g_\ell
    =
    \sigma(w_\ell^\top c_\ell + b_\ell),
\end{equation}
and then applied only to depth-token keys and values:
\begin{equation}
    \bar{K}^{\mathrm{vlm}}_{\ell,t}
    =
    \left(1-M_t+M_t g_\ell\right)K^{\mathrm{vlm}}_{\ell,t},
    \qquad
    \bar{V}^{\mathrm{vlm}}_{\ell,t}
    =
    \left(1-M_t+M_t g_\ell\right)V^{\mathrm{vlm}}_{\ell,t}.
\end{equation}
The gated \((\bar{K}^{\mathrm{vlm}}_{\ell},\bar{V}^{\mathrm{vlm}}_{\ell})\)
are then projected into the action expert as in \autoref{sec:posttraining}.
The gate is initialized with bias \(-4\), so fine-tuning begins close to the
standard action-conditioning path and learns how strongly each expert layer
should use the depth prefix.
